% Level 2, Week 7 -- Files: Saving Your Stuff (student workbook)
\chapterweek{7}{Files: Saving Your Stuff}{A program forgets everything when it closes, unless it writes things down in a file.}

\section{What you will learn}
\begin{itemize}[leftmargin=1.2em]
  \item How to \textbf{save} data in a text file, and how to \textbf{load} it again next time.
  \item The three ways to open a file: \texttt{"r"} read, \texttt{"w"} write, \texttt{"a"} append.
  \item How to clean up lines with \texttt{strip()} and cut them into pieces with \texttt{split()}.
  \item What to do when the file is missing, and what a game should never save about you.
\end{itemize}

\section{The big idea: programs forget, files remember}
When a program closes, all its variables are wiped. Your high score? Gone. A \textbf{file} is like a notebook in a drawer:
it stays on the computer after the program closes, so next time the program can read it again.

\begin{center}
%% alt: How a program remembers things. Your program's variables live inside the program and are wiped when it closes. Writing with mode "w" or "a" sends data to the file high_scores.txt, which holds Ada,120 and Kai,95. Reading with mode "r" brings it back. The file stays on the disk next to your program.
\begin{tikzpicture}[font=\small,>={Stealth[length=3mm]},
  box/.style={draw=black,thick,rounded corners=3pt,minimum width=3.2cm,minimum height=1.6cm,align=center}]
  \node[box,fill=kcsky!25] (prog) at (0,0) {\textbf{Your program}\\variables live here\\(wiped when it closes)};
  \node[box,fill=kcyellow!35] (file) at (7,0) {\textbf{high\_scores.txt}\\\texttt{Ada,120}\\\texttt{Kai,95}};
  \draw[->,very thick,kcblue] ([yshift=4mm]prog.east) -- node[above,text=black]{\textbf{write}: \texttt{"w"} or \texttt{"a"}} ([yshift=4mm]file.west);
  \draw[->,very thick,kcvermillion,dashed] ([yshift=-4mm]file.west) -- node[below,text=black]{\textbf{read}: \texttt{"r"}} ([yshift=-4mm]prog.east);
  \node[below=1mm of file,font=\footnotesize] {stays on the disk, next to your program};
\end{tikzpicture}
\end{center}

\subsection{Writing to a file}
\begin{pycode}[Save one score]
with open("high_scores.txt", "a") as file:
    file.write("Ada,120\n")
\end{pycode}
\begin{itemize}[leftmargin=1.2em]
  \item \texttt{open(name, mode)} opens the file. \texttt{with} closes it again for you when the indented block ends.
  \item \texttt{\textbackslash n} means \enquote{new line}. Without it, everything ends up squashed on one line.
\end{itemize}

\begin{center}
\renewcommand{\arraystretch}{1.3}
\begin{tabular}{@{}>{\ttfamily}l p{4.2cm} p{6.4cm}@{}}
\toprule
\textrm{\textbf{Mode}} & \textbf{Means} & \textbf{Like a notebook\ldots} \\ \midrule
"r" & read & Just read it. Crashes if the notebook is not there. \\
"a" & append (add to the end) & Write on the next empty line. Makes a new notebook if needed. \\
"w" & write (replace) & Rip out every page, then write. \textbf{Careful: old data is gone!} \\
\bottomrule
\end{tabular}
\end{center}

\subsection{Reading a file, line by line}
\begin{pycode}[Load the scores back]
scores = []
with open("high_scores.txt", "r") as file:
    for line in file:
        line = line.strip()            # "Ada,120\n" -> "Ada,120"
        if line != "":                 # skip empty lines
            parts = line.split(",")    # "Ada,120"   -> ["Ada", "120"]
            scores.append([parts[0], int(parts[1])])
print(scores)                          # [['Ada', 120], ['Kai', 95]]
\end{pycode}
Everything you read from a file is \textbf{text} (a \texttt{str}), just like \texttt{input()}. Use \texttt{int()} to turn \texttt{"120"} into the number \texttt{120}.

\subsection{Bug hunt: the file is not there!}
The very first time a game runs, there is no score file yet. Reading it gives:
\begin{outputcode}
FileNotFoundError: [Errno 2] No such file or directory: 'high_scores.txt'
\end{outputcode}
Check the spelling of the name first. If the file really is new, catch the error with \texttt{try} and \texttt{except}:
\begin{pycode}
try:
    with open("high_scores.txt", "r") as file:
        print(file.read())
except FileNotFoundError:
    print("No scores yet - you will be the first!")
\end{pycode}
\texttt{try} means \enquote{have a go}. \texttt{except FileNotFoundError} means \enquote{if the file is missing, do this instead of crashing}.

\begin{exercise}[Save High Scores]
Folder: \texttt{exercises/week07/ex-save-high-scores}. Open \texttt{starter.py}, press \menu{Run} twice. Does it remember your first game?
\begin{enumerate}[leftmargin=1.4em]
  \item TODO 1: in \texttt{save\_score()}, open the file with \texttt{"a"} and write one line \texttt{nickname,score}.
  \item TODO 2: in \texttt{load\_scores()}, strip, skip empty lines, split at the comma and append \texttt{[nickname, score]}.
  \item TODO 3: if the file is missing, make an empty one with \texttt{"w"}.
  \item Play three times. Open \texttt{high\_scores.txt} (Finder: double-click it). Copy the first line here:
  \answerlines{1}
  \item Check your work with \texttt{check\_my\_work.py}.
\end{enumerate}
\end{exercise}

\begin{exercise}[Word Guess (choose this one or 7.3)]
Folder: \texttt{exercises/week07/ex-word-guess}. The secret words live in \texttt{words.txt}, one per line.
\begin{enumerate}[leftmargin=1.4em]
  \item TODO 1: \texttt{load\_words()} reads every line, strips it, makes it lowercase with \texttt{lower()} and skips empty lines.
  \item TODO 2: \texttt{show\_progress()} shows guessed letters and \texttt{\_} for the rest, like \texttt{p \_ \_ e \_}.
  \item Play a round with your partner, then add three new words to \texttt{words.txt}. My new words:
  \answerlines{1}
\end{enumerate}
\end{exercise}

\begin{exercise}[Dot Art from a File (choose this one or 7.2)]
Folder: \texttt{exercises/week07/ex-dot-art-file}. Each line of \texttt{my\_art.txt} is one dot: \texttt{x,y,size,color}.
\begin{enumerate}[leftmargin=1.4em]
  \item TODO 1: \texttt{parse\_line()} turns \texttt{"-150,-60,40,\#009E73"} into \texttt{[-150, -60, 40, "\#009E73"]}.
  \item TODO 2: \texttt{load\_dots()} reads every line that is not empty. Run it: a caterpillar under a sun!
  \item TODO 3: \texttt{save\_dots()} writes every dot back, one per line. Which mode is right here, and why?
  \answerlines{1}
  \item Change \texttt{my\_art.txt} to draw something new. Colors: \swatch{kcorange}~\texttt{\#E69F00} \swatch{kcsky}~\texttt{\#56B4E9} \swatch{kcgreen}~\texttt{\#009E73}
        \swatch{kcyellow}~\texttt{\#F0E442} \swatch{kcblue}~\texttt{\#0072B2} \swatch{kcvermillion}~\texttt{\#D55E00} \swatch{kcpurple}~\texttt{\#CC79A7}.
\end{enumerate}
\end{exercise}

\begin{tip}[What should a game save?]
A file stays on the computer, and anyone who uses the computer can open it. Save only what the program needs:
nicknames, scores, words, drawings. Never save real names, addresses, schools, birthdays or passwords.
\end{tip}

\begin{vocab}
  \vocabword{file}{Data saved on the computer's disk. It is still there after the program closes.}
  \vocabword{open / with}{\texttt{open} gets a file ready to use; \texttt{with} closes it again automatically.}
  \vocabword{read / write / append}{The modes \texttt{"r"}, \texttt{"w"} and \texttt{"a"}: look, replace everything, or add to the end.}
  \vocabword{strip()}{Removes spaces and the invisible new-line \texttt{\textbackslash n} from both ends of a string.}
  \vocabword{split(",")}{Cuts a string into a list at every comma: \texttt{"a,b"} becomes \texttt{["a", "b"]}.}
  \vocabword{FileNotFoundError}{The error you get when you try to read a file that does not exist.}
  \vocabword{data privacy}{Deciding what information about people is saved, who can see it, and keeping it safe.}
\end{vocab}

\begin{reflection}
What is the difference between a variable and a file? \answerlines{2}
My partner's game saved \enquote{only the newest score}. What do you think the bug was? \answerlines{1}
Name two things a game should never save about you. \answerlines{1}
\end{reflection}
